% Part: first-order-logic
% Chapter: sequent-calculus
% Section: soundness-identity

\documentclass[../../../include/open-logic-section]{subfiles}

\begin{document}

\olfileid{fol}{seq}{sid}

\olsection{Correctheid met \usetoken{S}{identity}}

\begin{prop}
$\Log{LK}$ blijft correct wanneer we beginsequenten en regels voor
identiteit, dus gelijkheid, toevoegen.
\end{prop}

\begin{proof}
Beginsequenten van de vorm ${} \Sequent \eq[t][t]$ zijn geldig. De
reden is dat voor elke !!{structure}~$\Struct M$ geldt
$\Sat{M}{\eq[t][t]}$. We nemen hierbij aan dat de term $t$ gesloten is:
hij bevat geen variabelen. Een toekenning van objecten aan variabelen
maakt daarom geen verschil.

Stel nu dat de laatste afleidingsstap in !!a{derivation} een $=$-regel
is. De premisse is $\eq[t_1][t_2], \Gamma \Sequent \Delta, !A(t_1)$;
de conclusie is $\eq[t_1][t_2], \Gamma \Sequent \Delta, !A(t_2)$.
Beschouw !!a{structure}~$\Struct M$. Om de conclusie geldig te noemen,
moeten we het volgende aantonen. Als $\Sat{M}{\eq[t_1][t_2]}$ en
$\Sat{M}{\Gamma}$, dan geldt ofwel $\Sat{M}{!C}$ voor een $!C \in
\Delta$, ofwel $\Sat{M}{!A(t_2)}$.

Volgens de inductiehypothese is de premisse geldig. Als
$\Sat{M}{\eq[t_1][t_2]}$ en $\Sat{M}{\Gamma}$, zijn er dus twee
mogelijkheden: (a) voor een $!C \in \Delta$ geldt $\Sat{M}{!C}$, of
(b) $\Sat{M}{!A(t_1)}$. In geval (a) zijn we meteen klaar. Neem nu
geval (b). Laat $s$ een toekenning zijn waarvoor
$s(x) = \Value{t_1}{M}$. Volgens
\olref[syn][ass]{prop:sentence-sat-true} geldt
$\Sat{M}{!A(t_1)}[s]$. Omdat $\varAssign{s}{s}{x}$, kunnen we
\olref[syn][ext]{prop:ext-formulas} dat $\Sat{M}{!A(x)}[s]$. Uit
$\Sat{M}{\eq[t_1][t_2]}$ volgt $\Value{t_1}{M} = \Value{t_2}{M}$, en
daarom $s(x) = \Value{t_2}{M}$. We passen nu
\olref[syn][ext]{prop:ext-formulas} nogmaals toe te passen, verkrijgen
we ook $\Sat{M}{!A(t_2)}[s]$. Ten slotte volgt volgens
\olref[syn][ass]{prop:sentence-sat-true} geldt derhalve
$\Sat{M}{!A(t_2)}$.
\end{proof}
\end{document}
